%==============================================================================
%  mammoviz-tda — Guía de uso del entorno de trabajo
%  Documento independiente de la memoria del TFG (tfg/main.tex)
%==============================================================================
\documentclass[11pt,a4paper]{article}

\usepackage[utf8]{inputenc}
\usepackage[T1]{fontenc}
\usepackage[spanish,es-tabla]{babel}
\usepackage{lmodern}

\usepackage[a4paper,top=2.5cm,bottom=2.5cm,left=2.8cm,right=2.8cm]{geometry}
\usepackage{setspace}
\onehalfspacing
\usepackage{emptypage}

\usepackage{graphicx}
\usepackage{float}
\usepackage{booktabs}
\usepackage{array}
\usepackage{caption}
\usepackage{xcolor}
\definecolor{bg}{rgb}{0.96,0.96,0.96}
\definecolor{accent}{rgb}{0.11,0.35,0.60}

\usepackage{listings}
\lstset{
  basicstyle=\ttfamily\small,
  backgroundcolor=\color{bg},
  keywordstyle=\color{blue!70!black}\bfseries,
  commentstyle=\color{green!40!black},
  stringstyle=\color{red!60!black},
  frame=single, breaklines=true, showstringspaces=false,
  captionpos=b, language=bash, columns=flexible,
}
\newcommand{\cmd}[1]{\lstinline[language=bash]|#1|}

\usepackage[hidelinks,colorlinks=true,linkcolor=accent,urlcolor=accent]{hyperref}
\usepackage{enumitem}
\usepackage{fancyhdr}
\pagestyle{fancy}
\fancyhf{}
\fancyhead[L]{mammoviz-tda — guía de uso}
\fancyhead[R]{\thepage}
\renewcommand{\headrulewidth}{0.4pt}

\title{\textbf{mammoviz-tda}\\[4pt]\large Guía de uso del entorno de trabajo}
\author{}
\date{Noviembre 2026}

\begin{document}
\maketitle
\thispagestyle{empty}

\begin{abstract}
\noindent Este documento es una guía práctica, independiente de la memoria
del TFG, para instalar el entorno, ejecutar el pipeline con datos reales de
CBIS-DDSM y trabajar de forma interactiva en un notebook. No reemplaza la
memoria (\texttt{tfg/main.pdf}): aquí solo se documenta \emph{cómo usar el
código}.
\end{abstract}

\tableofcontents

%==============================================================================
\section{Visión general}

El proyecto combina Análisis Topológico de Datos (TDA) y \textit{deep
learning} sobre mamografías reales de \href{https://www.cancerimagingarchive.net/collection/cbis-ddsm/}{CBIS-DDSM},
con el objetivo de estudiar la detección anticipada de cáncer de mama mediante
comparación longitudinal de exploraciones sucesivas (ver los issues del
milestone \texttt{F0'} en GitHub para el contexto completo de esta hipótesis).

\subsection{Estructura del repositorio}

\begin{center}
\begin{tabular}{@{}ll@{}}
\toprule
\textbf{Carpeta/fichero} & \textbf{Contenido} \\
\midrule
\texttt{src/} & Librería: carga de datos, preprocesado, TDA, modelos, evaluación \\
\texttt{scripts/} & Ejecutables de un solo comando (ver Sección~\ref{sec:comandos}) \\
\texttt{notebooks/} & Entorno de trabajo interactivo (ver Sección~\ref{sec:notebook}) \\
\texttt{tests/} & Tests automáticos (\texttt{pytest}) \\
\texttt{data/} & Datos descargados (NO versionado en git; hay que descargarlos) \\
\texttt{results/} & Salidas de los scripts (métricas, figuras) \\
\texttt{tfg/} & Memoria del TFG en LaTeX \\
\texttt{bib/} & PDFs de la bibliografía \\
\bottomrule
\end{tabular}
\end{center}

\textbf{Importante — qué NO persiste entre sesiones.} La carpeta \texttt{data/}
está en \texttt{.gitignore} a propósito (163\,GB en el dataset completo,
imposible de versionar). Cada vez que empiece de cero en un equipo nuevo, o si
la carpeta desaparece por limpieza/sincronización, hay que volver a ejecutar
los comandos de descarga de la Sección~\ref{sec:comandos}.

%==============================================================================
\section{Instalación}

\begin{lstlisting}
# 1. Entorno virtual
python -m venv .venv
.venv\Scripts\activate          # Windows
source .venv/bin/activate       # Linux/Mac

# 2. Dependencias
pip install -r requirements.txt

# 3. Verificar que todo funciona
pytest
\end{lstlisting}

\texttt{requirements.txt} incluye \texttt{gudhi}/\texttt{giotto-tda} (TDA),
\texttt{torch}/\texttt{torchvision} (CNN), \texttt{pydicom}/\texttt{requests}
(descarga de datos reales) y \texttt{pytest}. En Python~3.13+ \texttt{giotto-tda}
no publica ruedas: el código detecta automáticamente qué backend TDA hay
instalado (\texttt{src.topology.backend()}) y usa \texttt{gudhi} de respaldo.

Si \texttt{pytest} da 45 tests en verde (1 \textit{skip} esperado si no tienes
\texttt{giotto-tda}), el entorno está listo.

%==============================================================================
\section{Comandos, en el orden en que se usan}
\label{sec:comandos}

\subsection{1. Descargar los metadatos reales (CSV)}

Los 4 CSV de descripción de casos de CBIS-DDSM (con la etiqueta real
benigno/maligno) se descargan directamente, sin necesidad de registro:

\begin{lstlisting}
mkdir -p data/csv
curl -O https://www.cancerimagingarchive.net/wp-content/uploads/mass_case_description_train_set.csv
curl -O https://www.cancerimagingarchive.net/wp-content/uploads/mass_case_description_test_set.csv
curl -O https://www.cancerimagingarchive.net/wp-content/uploads/calc_case_description_train_set.csv
curl -O https://www.cancerimagingarchive.net/wp-content/uploads/calc_case_description_test_set.csv
# mover los 4 ficheros .csv descargados a data/csv/
\end{lstlisting}

\subsection{2. Descargar imágenes reales (piloto incremental)}

El dataset completo de imágenes son 163{,}5\,GB — demasiado para descargar de
golpe. Este comando descarga un piloto balanceado benigno/maligno
directamente contra la API REST de TCIA (no hace falta el NBIA Data
Retriever), y se puede volver a lanzar cuantas veces se quiera: nunca repite
descargas ya hechas, solo añade pacientes nuevos.

\begin{lstlisting}
python scripts/download_cbis_ddsm.py --add 15 15 --category mass
\end{lstlisting}

\texttt{--add 15 15} añade 15 pacientes benignos y 15 malignos \emph{nuevos}.
Tarda varios minutos (cada mamografía completa son decenas de MB). Se puede
relanzar en segundo plano y comprobar el progreso en el fichero de log.

\subsection{3. Pipeline TDA completo}

Un único comando: carga datos reales $\to$ partición por paciente (sin fuga
de datos) $\to$ preprocesado $\to$ homología persistente $\to$ clasificador
$\to$ métricas clínicas.

\begin{lstlisting}
python scripts/run_pipeline.py
python scripts/run_pipeline.py --classifier rf --n-bins 4   # valores por defecto
python scripts/run_pipeline.py --tune                       # ajusta hiperparametros por GridSearchCV
\end{lstlisting}

Guarda en \texttt{results/}: \texttt{confusion\_matrix.png},
\texttt{persistence\_diagram\_example.png} y \texttt{metrics.txt} (incluye el
AUC por validación cruzada, más fiable que el de un test pequeño).

\subsection{4. CNN de referencia (PyTorch)}

\begin{lstlisting}
python scripts/train_cnn.py --epochs 10
\end{lstlisting}

Entrena una DenseNet121 (base congelada, cabeza nueva) sobre las mismas
imágenes y partición que el pipeline TDA. Semilla fija: resultados
reproducibles entre ejecuciones.

\subsection{5. Comparar TDA vs. CNN vs. fusión}

\begin{lstlisting}
python scripts/fuse_and_evaluate.py
\end{lstlisting}

Entrena ambos modelos sobre los mismos casos y reporta accuracy/AUC (test y
validación cruzada) para las 3 vías: TDA sola, CNN sola, fusión de
características (\texttt{src.models.fuse\_features}).

\subsection{6. Validación de la hipótesis longitudinal (datos sintéticos)}

Mientras se gestiona el acceso a un dataset longitudinal real (CSAW-CC o
EMBED, ambos bajo solicitud formal — ver issue \texttt{\#92} en GitHub):

\begin{lstlisting}
python scripts/synthetic_longitudinal_validation.py
python scripts/synthetic_longitudinal_validation.py --sweep   # analisis de sensibilidad
\end{lstlisting}

%==============================================================================
\section{Entorno de trabajo interactivo: notebook}
\label{sec:notebook}

Para explorar, depurar o visualizar resultados celda a celda (en vez de
ejecutar un script de un tirón), usa:

\begin{center}
\texttt{notebooks/02\_real\_data\_pipeline.ipynb}
\end{center}

Reutiliza exactamente las mismas funciones que los scripts de la
Sección~\ref{sec:comandos} (\texttt{run\_pipeline.py},
\texttt{train\_cnn.py}, \texttt{synthetic\_longitudinal\_validation.py}), así
que nunca hay dos implementaciones del mismo cálculo que puedan divergir.
Las secciones del notebook son:

\begin{enumerate}[leftmargin=1.2cm]
  \item Imports y configuración.
  \item Carga de metadatos reales + partición por paciente.
  \item Visualización de un caso real (mamografía + máscara ROI).
  \item Preprocesado y diagrama de persistencia de ese mismo caso.
  \item Construcción del dataset completo (todos los casos ya descargados).
  \item Clasificador topológico + métricas clínicas.
  \item Validación por validación cruzada (más fiable que un test pequeño).
  \item CNN de referencia + fusión (opcional, más lento).
  \item Validación sintética de la hipótesis longitudinal.
\end{enumerate}

\begin{lstlisting}
jupyter notebook notebooks/02_real_data_pipeline.ipynb
\end{lstlisting}

El notebook original con datos sintéticos
(\texttt{notebooks/01\_pipeline\_demo.ipynb}) sigue disponible para quien
quiera ejecutar sin tener CBIS-DDSM descargado.

\subsection{Si quieres encadenar varios notebooks}

Para un flujo de trabajo más largo (p.\,ej. separar exploración de datos,
experimentos de TDA y experimentos de CNN en notebooks distintos), la forma
recomendada de \emph{no} duplicar código es que cada notebook nuevo importe
las mismas funciones compartidas, igual que hace
\texttt{02\_real\_data\_pipeline.ipynb}:

\begin{lstlisting}
import sys
from pathlib import Path
sys.path.insert(0, str(Path.cwd().parent))          # raiz del repo
sys.path.insert(0, str(Path.cwd().parent / "scripts"))

import run_pipeline as rp
import train_cnn as tc
import synthetic_longitudinal_validation as slv
\end{lstlisting}

Así cada notebook es una \emph{vista} distinta sobre el mismo pipeline, no una
copia del código.

%==============================================================================
\section{Problemas conocidos y cómo evitarlos}

\begin{description}
  \item[\texttt{data/} ha desaparecido.] Normal (está en \texttt{.gitignore}).
    Repite los pasos 1 y 2 de la Sección~\ref{sec:comandos}.
  \item[Un script falla con \texttt{SyntaxError} en un marcador de conflicto de Git.]
    El fichero quedó con un conflicto de Git sin
    resolver (pasó una vez con \texttt{scripts/run\_pipeline.py}, causado por
    una herramienta de backup automático que hace commits sin supervisión).
    Busca y resuelve el conflicto a mano, o pide ayuda para resolverlo.
    Verifica que no quede ninguno con:
    \begin{lstlisting}
grep -rln "^<<<<<<<" --include="*.py" .
    \end{lstlisting}
  \item[Backups automáticos revierten cambios recientes.] Se ha observado que
    algún commit automático (\texttt{"Backup automatico: ..."}) restaura
    versiones antiguas de ficheros concretos. Si algo que arreglaste
    "vuelve" a su estado anterior, no es un error tuyo: compruébalo con
    \texttt{git log --oneline -- <fichero>} y vuelve a aplicar el cambio.
  \item[AUC muy distinto entre ejecuciones.] Con pilotos pequeños (pocas
    decenas de pacientes), el AUC sobre un test de pocos casos es ruidoso.
    Usa siempre el AUC por validación cruzada que reportan los scripts, no
    solo el del test.
\end{description}

\end{document}
